\origpage{377}{N. Robertson, P.D. Seymour / Journal of Combinatorial Theory, Series B 90 (2004) 325--385}
\noindent say; and let $(C, D)$ be the $(\mathscr{T}, \mathscr{T}')$-distinction. Let $F$ be a $\mathscr{T}$-enclave around $\mathscr{T}'$. Then $|E(F)| \le 4\rho(W) + 4$. If $(A, B)$ is the hub sector at $W$, then some edge of $A$ is included in $ins(F)$, and so by 8.6, either $ins(H(W)) \subseteq ins(F)$ or $ins(F) \subseteq ins(H(W))$. Now since $|E(F)| \le 4\rho(W) + 4$ and $H(W)$ is a $W$-hub, we deduce that if $ins(H(W)) \subseteq ins(F)$ then $F = H(W)$. Thus if $(A, B)$ is the hub sector then $ins(F) \subseteq ins(H(W))$ and so $C \subseteq A$ and $B \subseteq D$ as required. On the other hand, if $(A, B)$ is an $r^*$-sector for some $r^* \in (V(\Gamma^*) \setminus V(\Gamma)) \cap bd(W)$, then by 13.6 $ins(F) \subseteq \Delta_{r^*}$, and again $C \subseteq A$ and $B \subseteq D$. Thus $\mathscr{L}$ $(4n \cdot 5^{2m} + 3)$-isolates $\mathscr{T}$. This proves (2).

Let $\mathscr{C}$ be the set of all $(\Gamma, \phi, \mathscr{L})$ as above such that $(\Gamma, \phi, \lambda, \mathscr{T}) \in \mathscr{D}'$ for some $\lambda, \mathscr{T}$. It remains to show that $\mathscr{C}$ is well-behaved. Thus, let $\Omega$ be a well-quasi-order, and let $(\Gamma_i, \phi_i, \mathscr{L}_i)$ $(i = 1, 2)$ be a countable sequence of members of $\mathscr{C}$. For each $i \ge 1$, let $\xi_i : \mathscr{L}_i \to E(\Omega)$ be some function. We must show that there exist $j > i \ge 1$ and an outline $\tau : \mathscr{L}_i \to \mathscr{L}_j$ such that $\xi_i(A) \le \xi_j(\tau(A))$ for all $A \in \mathscr{L}_i$. Let $(\Gamma_i, \phi_i, \lambda_i, \mathscr{T}_i) \in \mathscr{D}'$ for each $i \ge 1$.

For each $W$, if $|V(H(W))| \ne 1$, let $A_i(W)$ be the hub sector of $(\Gamma_i, \phi_i, \lambda_i, \mathscr{T}_i)$ at $W$. There is an infinite set $I \subseteq \{1, 2, \dots\}$ such that for all $i, j \in I$ with $j > i$, $\xi_i(A_i(W)) \le \xi_j(A_j(W))$ because $\Omega$ is a well-quasi-order. To simplify notation, let us replace our initial sequence by this subsequence. We may therefore assume that

\hspace*{12pt}(3) \textit{For each $W \in \mathscr{W}$ with $|V(H(W))| \ne 1$, $\xi_1(A_1(W)) \le \xi_2(A_2(W)) \le \dots$}

For each $r^* \in (V(\Gamma_i^*) \setminus V(\Gamma_i)) \cap bd(W)$, let $A_i(r^*)$ be the sector of $(\Gamma_i, \phi_i, \lambda_i, \mathscr{T}_i)$ at $r^*$. We may similarly assume that

\hspace*{12pt}(4) \textit{For each $W \in \mathscr{W}$ and each $r^* \in S(W)$, $\xi_1(A_1(r^*)) \le \xi_2(A_2(r^*)) \le \dots$}

{\sloppy Let $R$ be the well-quasi-order with element set all \del{triples} \fix{quintuples} $(\pi, \pi_1, \pi_2\fix{, h_1, h_2})$, where $\pi$ is a march in $\Sigma_\chi$ with $|\bar{\pi}| \le 4n \cdot 5^{2m} + 2$, and $\pi_1, \pi_2$ are marches with $\bar{\pi}_1, \bar{\pi}_2 \subseteq \bar{\pi}$\fix{, and $h_1, h_2$ are integers with $0 \le h_1, h_2 \le 4n \cdot 5^{2m} + 2$}. We order $R$ by isomorphism; that is, $(\pi, \pi_1, \pi_2\fix{, h_1, h_2}) \le (\pi', \pi'_1, \pi'_2\fix{, h'_1, h'_2})$ if \fix{$h_1 = h'_1$, $h_2 = h'_2$, and} there is a bijection of $\bar{\pi}$ to $\bar{\pi}'$ mapping $\pi$ to $\pi'$, $\pi_1$ to $\pi'_1$ and $\pi_2$ to $\pi'_2$. For each $i \ge 1$, each $W \in \mathscr{W}$, and each regular $r^* \in (V(\Gamma_i^*) \setminus V(\Gamma_i)) \cap bd(W)$, let $e_{r^*}$ be the component of $bd(W) \setminus \{v_1, v_2\}$ containing $r^*$, where $v_1, v_2$ are the neighbours of $r^*$ in $\Gamma_i^* \cap bd(W)$. Let
\begin{equation*}
\phi'_i(e_{r^*}) = (\xi(A_i(r^*)), (\pi(A_i(r^*)), \mu_i(v_1), \mu_i(v_2)\fix{, h_i(v_1), h_i(v_2)}),
\end{equation*}
where $\mu_i(v_1), \mu_i(v_2)$ are the marches given by enumerating the vertices of $V(S_{v_1}) \cap V(\Gamma_i)$, $V(S_{v_2}) \cap V(\Gamma_i)$ (respectively) in order, starting with $v_1$ and $v_2$\fix{; and for $v \in V(\Gamma_i) \cap bd(W)$, with $\mu_i(v)$ defined in the same way, the \textit{hub index} $h_i(v)$ is $h$ if the last term of $\mu_i(v)$ is the $h$th term of $\pi(H(W))$, and $0$ if the last term of $\mu_i(v)$ is not a term of $\pi(H(W))$}. Then $\phi'_i(e_{r^*}) \in \Omega \times R$.\fixnote{fix:14.1-hub-index}{The hub indices are not in the paper. They are used in claim (10) of the construction of $\sigma$ that replaces the paper's, on page 379. Since $|\bar{\pi}(H(W))| \le 2\rho(W) + 2 \le 4n \cdot 5^{2m} + 2$, the hub indices are within the stated bounds; and $R$ is still a well-quasi-order, since it has only finitely many classes of mutually equivalent elements.}\par}

Let $\Sigma'$ be obtained from $\Sigma_\chi$ by deleting $\Sigma_\chi \cap (W \setminus bd(W))$ for each $W \in \mathscr{W}$. Then $\Sigma'$ is a surface since the $W$'s are disjoint and each $bd(W)$ is disjoint from $bd(\Sigma_\chi)$. Moreover, $bd(\Sigma') \cap bd(\Sigma_\chi) = \emptyset$, and $\hat{\Sigma} \cong \hat{\Sigma}_\chi$. Let $\Phi'$ be the frame with $U(\Phi') = bd(\Sigma')$, $V(\Phi') = \bigcup (N(W) : W \in \mathscr{W})$, where the edges of $\Phi'$ are directed arbitrarily, and an edge of $\Phi'$ is short if it contains a vertex of $S(W)$ for some $W \in \mathscr{W}$, and long otherwise. Let $\chi'$ be the colour scheme where $\Sigma_{\chi'} = \Sigma'$, $\Phi_{\chi'} = \Phi'$, $\Omega_{\chi'}(k)$ is an ideal of
